\section{Blow-ups and the dependence on the compactification}
\label{sec:compactifications}

Here \emph{good} means that every connected stratum is aspherical and that,
after choosing a path from a small normal torus to a base point of the relevant
component of $U=X\setminus D$, its meridian homomorphism is injective.
A different path conjugates this homomorphism and does not change injectivity.
The assertions about meridians below concern actual homomorphisms after these
choices, not merely their outer conjugacy classes.  We first work with complex
manifolds; the local statements descend through an \'etale atlas of a smooth
analytic Deligne--Mumford stack.  In the surface classification, \emph{completion}
means a proper algebraic completion with the inclusion of $U$ fixed.

\begin{theorem}[\status{Proved}: the local blow-up calculation]\label{P2a}
Let $Z\subset D$ be a smooth closed centre having normal crossings with $D$,
and let $\beta:(X',D')\to(X,D)$ be its blow-up, with reduced total transform.
At a general point of a component of $Z$, choose local coordinates in which
\[
 D=\{x_1\cdots x_s=0\},\qquad
 Z=\{x_1=\cdots=x_s=y_1=\cdots=y_t=0\}.
\]
Here $s\geq1$, and $\operatorname{codim}_X Z=s+t$.  The exceptional meridian
has image $m_E=m_1\cdots m_s$ in the local group $\Z^s$ and therefore in
$\pi_1(U)$.  The factors commute in the specified local group.
If $t>0$, the exceptional divisor and all $s$ strict transforms meet.
At every point of their open common intersection the new lattice map factors as
\[
 \Z^{s+1}\longrightarrow\Z^s\longrightarrow\pi_1(U),\qquad
 (a_1,\ldots,a_s,b)\longmapsto(a_1+b,\ldots,a_s+b).
\]
In particular, $(-1,\ldots,-1,1)$ is in its kernel.  The same nonzero primitive
vector gives a kernel in the profinite lattice $\widehat\Z(1)^{s+1}$.
If $t=0$, each maximal new cone has generators
\[
 e_1,\ldots,\widehat e_j,\ldots,e_s,
 e_E=e_1+\cdots+e_s;
\]
its change-of-basis matrix has determinant $\pm1$.
Consequently a boundary point blown up on one branch of a surface produces
failure at the unique new double point.  The exceptional smooth stratum has
meridian $m_1$ and need not itself fail injectivity.
\end{theorem}
\begin{proof}
In the blow-up chart where the exceptional homogeneous coordinate $y_1$ is
nonzero, write $x_i=y_1u_i$.  The reduced boundary is
$y_1u_1\cdots u_s=0$.  Rotating $u_i$ once rotates $x_i$ once, whereas rotating
$y_1$ once rotates every $x_i$ once.  This proves both the displayed matrix
and its kernel.  This chart exists precisely because $t>0$.  At a point of its
common intersection the $s+1$ local normal circles are independent before
mapping to the old complement.  Completing the integer matrix leaves the
primitive kernel vector nonzero; no residual-finiteness assumption on
$\pi_1(U)$ is needed for this failure.

If $t=0$, every homogeneous point of $\PP^{s-1}$ has some $x_j$ coordinate
nonzero, so the exceptional divisor cannot meet all $s$ strict transforms.
The chart $x_i=x_ju_i$ for $i\ne j$ gives exactly the stated basis.  Its
unimodularity follows by subtracting the other columns from the last column.
For $s=1,t=1$ this is the equality of the exceptional and old meridians at
$E\cap\widetilde D$.  It says nothing about the order of their common image;
thus it identifies exactly the obligatory new rank-two failure, rather than
asserting failure everywhere on $E$.
\end{proof}

This is the precise content supported by cote 149, folio 129, scan 130,
including the qualification ``en g\'en\'eral'' in the margin.  The statement
``$\operatorname{card}I(x)\ne1$'' is a warning about the newly created
intersection, not an assertion that every multiple boundary point is bad.
For example every toric boundary point of a smooth toric variety is good.

\begin{theorem}[\status{Proved}: functoriality, ends, and preservation]\label{P2b}
In the situation of Theorem~\ref{P2a}, the blow-up induces a functor
\[
 \Exit(X',D')\longrightarrow\Exit(X,D)
\]
by postcomposition of exit paths, and a functor on the categories of crossing
arrows with target in $U$.  Its restriction to $\Pi_\infty(U)$ is the identity.
The germs of punctured neighbourhoods of $D'$ and $D$ are isomorphic as
pro-spaces over $U$.  For compact normal-crossings pairs, their links at infinity
are consequently homotopy equivalent.

The pair $(X',D')$ is good if and only if $(X,D)$ is good when the centre is
locally a union of closures of strata.  Every genuinely nonstratal component
of the centre destroys the meridian condition.  Thus the literal phrase
``each centre which is not a closure of one stratum'' needs the correction
``each centre which is not locally a union of stratum closures'': disjoint
unions of permissible centres give no counterexample to preservation.
\end{theorem}
\begin{proof}
Use the stratifications refined by the connected components of branch flags.
Every new stratum maps into one old stratum.  The relation between its incident
strata is preserved by $\beta$, as follows directly from its monomial coordinate
expressions.  Postcomposition therefore preserves exit simplices and their
composition.  Applying the arrow-category construction and restricting the
second endpoint to $U$ gives the second functor.  In normal charts its map of
crossing links is the torus map of Theorem~\ref{P2a}, together with the map of
flag bases induced by $\beta$.

Because $\beta$ is proper, if $N'$ is a neighbourhood of $D'$, the closed image
$\beta(X'\setminus N')$ is disjoint from $D$.  Its complement $N$ satisfies
$\beta^{-1}N\subset N'$.  Conversely, inverse images of neighbourhoods of $D$
are neighbourhoods of $D'$.  These two systems are cofinal, and their punctured
members are the same open subsets of $U$, since $\beta$ is an isomorphism there.
This proves the pro-space assertion without an assumption of compactness.
For compact pairs a collar of the boundary of the real oriented blow-up gives
a cofinal system of punctured neighbourhoods all homotopy equivalent to that
boundary, proving the assertion about an actual link.

For a stratal centre, old and new local meridian lattices are identified by
unimodular bases.  Injectivity therefore transfers in both directions.  Over
an old stratum $S$ contained in the centre, the new strata are bundles over $S$
with fibres algebraic tori.  Their structure group may permute the torus
coordinates; after a finite flag cover it is the usual torus-bundle structure.
If $S$ is aspherical, the homotopy exact sequence of such a bundle makes its
total space aspherical.  Old strata not contained in the centre are unchanged.
This proves forward preservation of asphericity.

Conversely, suppose the new pair is good.  For each old $S$ use the open
exceptional stratum whose fibre is $(\C^*)^{s-1}$.  If $Q\to S$ is its
flag-covered torus bundle, goodness of $Q$ and the homotopy exact sequence give
$\pi_j(S)=0$ for $j\geq3$.  Old meridian injectivity has already been recovered
from the unimodular maximal charts.  Therefore the connecting homomorphism
$\pi_2(S)\to\Z^{I(S)}$ of the full old normal torus bundle is zero: its image
dies in that bundle and hence in $U$, where the lattice injects.  The torus
bundle $Q$ is the quotient of the relevant punctured normal factors by their
diagonal $\C^*$, so naturality of the homotopy exact sequences identifies its
connecting homomorphism with the projection of the preceding zero map to
$\Z^{s-1}$.  Since $\pi_2(Q)=0$, exactness now gives $\pi_2(S)=0$.
Finite flag covers do not change homotopy groups in degrees at least two.
Finally, if $t>0$ at a general point of any centre component,
Theorem~\ref{P2a} supplies the obligatory kernel, independently of the old pair.
All these arguments are invariant under the finite stabilizer actions in an
\'etale chart, and hence apply to smooth analytic DM stacks with representable
blow-up maps.
\end{proof}

\section{Minimal completions and elementary fibrations}

\subsection{A classification for algebraic surfaces}

For a proper normal-crossings completion $(Y,B)$ of a smooth algebraic surface,
form the finite graph whose vertices are the irreducible components of the
normalization of $B$, labelled by their geometric genera and self-intersections;
a self-node contributes a loop and hence two to the valence.  Its
\emph{core} is obtained by repeatedly deleting genus-zero vertices of valence
at most one.  Vertices of positive genus are never deleted.  A core pair is a
pair for which nothing is deleted.  This definition concerns the graph; deleting
a vertex graph-theoretically does not assert that the corresponding curve can
be contracted on the surface.

\begin{lemma}[\status{Proved}: core pairs and common corner refinements]\label{P3a-core}
Every good proper algebraic surface completion is a core pair.  Any two core
normal-crossings completions of the same marked algebraic surface have a common
refinement obtained from each by blowing up double points of the boundary.
Self-nodes are allowed among those double points.  In particular all core
completions lie in one corner birational class, whether or not their meridian
maps are injective.
\end{lemma}
\begin{proof}
An isolated rational boundary component is a stratum with nonzero $\pi_2$.
A rational component with just one boundary puncture has open stratum $\C$.
At that puncture the relative normal circle in this component is nullhomotopic.
Its image in the link of the open stratum is therefore an integral combination
of the normal meridian of the component.  The two local boundary meridians
cannot be independent.  Thus goodness excludes both rational valences zero
and one, and the entire boundary graph is its core.

Resolve the birational map between two completions by a smooth common
resolution $W$, with boundary its reduced inverse image.  Proper birational
morphisms between smooth algebraic surfaces factor into point blow-ups
\citep[Lemma 54.17.1]{StacksSurfaceFactorization2026}; resolving the graph supplies
such a common resolution.  All centres lie on the boundary, since the map is
the identity on $U$.  A blow-up of a node subdivides an edge, whereas a blow-up
of a smooth boundary point attaches a rational leaf.  Starting from a core
graph, its graph-theoretic complement in the new core is therefore a disjoint
union of rational trees, each attached to the core at one point.  Every
component outside the core of $W$ is exceptional for either projection:
strict transforms of components of either core pair belong to the core of $W$.

If the complement of the core is nonempty, choose the last point blow-up,
in a factorization over the first completion, that creates a component outside
the final core.  No later blow-up touches this component, since that would
create a later noncore component.  It is a smooth rational $(-1)$-curve of
valence one or two, so its contraction is smooth and preserves normal crossings.
It is exceptional for both projections, so both factor through this contraction
(the factorization property is also the content of Stacks, Lemma 54.16.1).
Repeat until the common resolution has no noncore components.  Neither
remaining point-blow-up factorization can contain a smooth-boundary centre:
such a blow-up would create a noncore tree that subsequent blow-ups cannot
remove.  Both factorizations therefore consist of corner blow-ups.
\end{proof}

\begin{theorem}[\status{Proved}: an intrinsic birational-end test for existence]\label{P3a-existence}
Let $U$ be a smooth connected algebraic surface admitting a proper smooth
normal-crossings completion.  Start with any such completion $(Y,B)$.
If some connected component of the boundary graph has empty core, declare the
core reduction unsuccessful.  Otherwise, contract smooth rational $(-1)$-curves
of valence one or two lying outside the graph core for as long as this is
possible.  If a noncore component remains, declare it unsuccessful.  If none
remains, denote the resulting core completion by $(Y_c,B_c)$.

Then the following criterion is independent of the initial completion and of
the permitted contraction choices:
\[
 \mathcal C(U)\ne\varnothing
 \quad\Longleftrightarrow\quad
 \begin{gathered}
 \text{core reduction succeeds, } U\text{ is aspherical, and}\[-2pt]
 \Z^{I(x)}\longrightarrow\pi_1(U)\text{ is injective at every stratum of }B_c.
 \end{gathered}                                      \tag{SC}
\]
Only finitely many meridian maps, one per connected stratum with its branch
flags, occur in this test.  If it succeeds, $\mathcal C(U)$ is exactly the
corner birational class of $(Y_c,B_c)$.  Thus the test uses one arbitrary end
presentation of $U$ and eliminates the existential search through other
completions; it is not an assertion that general finitely presented-group
injectivity or asphericity is algorithmically decidable.
\end{theorem}
\begin{proof}
Suppose first that a good completion $X$ exists.  In a common resolution of
$X$ and $Y$, all components outside the core are exceptional trees produced
by point blow-ups above smooth points of a boundary divisor.  They can be
contracted back to smooth points.  Follow a factorization of the morphism to
$Y$ into normal-crossings point blow-downs.  Contracting a noncore $(-1)$-curve
partially contracts one such exceptional tree and leaves a tree which still
contracts to a smooth point.  Contracting a core curve can only contract a
smooth rational bivalent vertex.  It cannot touch an attached noncore tree:
a core vertex carrying such a tree has total valence at least three, and its
contraction would create more than two local branches, contrary to normal
crossings.  Thus the property that each noncore tree contracts to a smooth
point persists all the way to $Y$.  Each boundary component has nonempty core,
and the stated reduction succeeds.  The same argument shows that any allowed
choice of a $(-1)$-curve leaves contractible exceptional trees, so a choice
cannot obstruct a successful reduction.

Its output is a core completion.  By Lemma~\ref{P3a-core} it has a common corner
refinement with $X$, and Theorem~\ref{P2b} transfers goodness to it.  Conversely,
if reduction succeeds, all positive-genus boundary strata are aspherical,
and every rational boundary stratum is a projective line minus at least two
points.  Zero-dimensional strata are points.  Hence the last two conditions
in (SC) say exactly that this particular completion is good.  The same common
corner-refinement lemma proves independence of the successful output and
identifies all other good completions.  If the graph core is empty in a boundary
component, point blow-ups only subdivide its tree or attach trees and cannot
produce a core.  A final isolated exceptional component cannot be contracted
while retaining a divisorial boundary and the same open $U$.  This also
justifies the first failure clause in the algorithm.
\end{proof}

\begin{theorem}[\status{Proved}: existence and uniqueness of minimal good members]\label{P3a-minimal}
If $\mathcal C(U)$ is nonempty, it has a member minimal under normal-crossings
boundary contractions.  Such minimal members need not be unique.
For a minimal member $(X,D)$, uniqueness as a marked completion is equivalent
to the following numerical conditions:
\begin{enumerate}[label=(\roman*)]
\item every smooth rational bivalent boundary component has self-intersection
at most $-2$;
\item every rational boundary component which is bivalent because it has exactly
one self-node and no other boundary intersection has self-intersection at most $2$.
\end{enumerate}
If the boundary is simple normal crossings, only condition (i) is needed.
Minimality here concerns boundary contractions preserving $U$ and normal
crossings; it is different from minimality of the underlying compact surface.
\end{theorem}
\begin{proof}
Every permitted contraction removes a boundary component, so repeated
contraction of smooth rational bivalent $(-1)$-components terminates.
Goodness persists by Theorem~\ref{P2b}.  These are all possible contractions in
a good pair: a rational tip is absent, a positive-genus curve is not exceptional
for a smooth surface contraction, and three surviving branches would violate
normal crossings.

Suppose (i)--(ii) hold and compare another minimal member using a common corner
refinement.  If any original component of $D$ disappears, consider the first
such contraction in a factorization to the other member.  A positive-genus
component cannot disappear.  A rational component with valence at least three
still has at least three branches, because no other original component has yet
been removed.  A smooth bivalent component has self-intersection at most its
original self-intersection: all previous contractions have removed only curves
exceptional over $X$, so the intermediate surface still maps to $X$ by point
blow-ups.  Its self-intersection is therefore at most $-2$ and it cannot be the
next $(-1)$-curve.  For a singly nodal bivalent component, making its strict
transform smooth requires first blowing up the node, which lowers its
self-intersection by $4$; it is consequently at most $-2$ whenever it is smooth.
Contracting the exceptional curve back restores the node and does not make the
original component contractible.  Thus no original component disappears.
The other completion dominates $X$ by corner blow-ups, and minimality forces
this morphism to be an isomorphism.

Conversely, let a smooth bivalent component have square $a\geq0$.  Blow up one
of its corners $a+1$ times along its successive strict transform.  Its square
is now $-1$; contract it and then contract any remaining admissible
$(-1)$-curves to reach a minimal member.  The original divisor valuation has
disappeared and cannot reappear under contractions, so this member is different
as a completion marked by $U$.  If a singly nodal bivalent component has square
$a\geq3$, first blow up its node; its smooth strict transform has square
$a-4\geq-1$.  Contract it immediately if the square is $-1$, or apply the
preceding construction otherwise.  Again its valuation disappears.
\end{proof}

\begin{example}[\status{Proved}: the six requested surface classes]\label{P3a-examples}
For algebraic completions marked by their specified open surface, the answers are
as follows.
\begin{center}
\begin{tabular}{p{.26\textwidth}p{.65\textwidth}}
\toprule
$U$ & $\mathcal C(U)$ and minimal members\\
\midrule
$\C^2$ & Empty.\\
$(\C^*)^2$ & Exactly the smooth complete toric surfaces with this marked torus.
Minimal members include $\PP^2$ and the Hirzebruch surfaces $\mathbb F_n$
for $n=0$ or $n\geq2$; these are all underlying minimal toric surface types.\\
$\C^*\times(\PP^1\setminus\{0,1,\infty\})$ & Corner refinements of ruled
completions with zero and infinity sections and the three missing fibres.
Minimal marked completions are indexed by line bundles on $\PP^1$ together
with a trivialization over the three-punctured line, equivalently divisors
supported at $0,1,\infty$ (an integer triple).\\
$C_1\times C_2$, both $C_i$ affine hyperbolic curves & All iterated corner
blow-ups of $\overline C_1\times\overline C_2$.  This product is the unique
minimal good completion.\\
$\PP^2\setminus Q$, $Q$ a smooth conic & Empty.\\
$\PP^2\setminus N$, $N$ a nodal cubic & Empty.\\
\bottomrule
\end{tabular}
\end{center}
\end{example}
\begin{proof}
A noncompact surface with finite fundamental group has no good completion:
its nonempty boundary supplies a lattice $\Z$ that cannot inject into that
finite group.  This settles $\C^2$ immediately.  For a smooth conic, identify
$\PP^2$ with $\operatorname{Sym}^2\PP^1$ and the conic with the diagonal locus.
The ordered configuration space of two distinct points of $S^2$ is simply
connected: projection to the first point has fibre $S^2\setminus\{p\}\simeq\R^2$.
Its free interchange quotient is the conic complement, so its fundamental group
is $\Z/2$.

For the nodal cubic, use $y^2=x^2(x+1)$ in an affine chart.  Projection to $x$
has one simple branch at $x=-1$ and one transverse node at $x=0$.
Van Kampen starts with the two meridians $a,b$ of a generic vertical line;
the simple branch imposes $a=b$, and the node imposes $[a,b]=1$.
Thus the affine complement has group $\Z$, generated by the irreducible-curve
meridian.  Adding the line at infinity kills its meridian, whose linking number
with a degree-three curve is three.  It imposes $a^3=1$, giving $\Z/3$.
The braid monodromies used here can be read directly from the roots
$\pm x\sqrt{x+1}$: around $-1$ they interchange once, and around $0$ they make
a full interchange squared.  There are no other finite critical values.
The finite presentations have been checked in
\texttt{code/compactifications\_checks.py}.

For the torus, all strata of a toric completion are tori and every local cone
is regular, so its primitive rays extend to a basis of $\Z^2=\pi_1(U)$.
A common corner refinement with a toric completion is toric; its corner
contractions remain toric.  Lemma~\ref{P3a-core} therefore proves that there
are no additional good algebraic completions.  The minimal regular complete
fans in dimension two give $\PP^2$ and $\mathbb F_n$, with $\mathbb F_1$
contracting to $\PP^2$.  Equivalently one can repeatedly remove a ray equal
to the sum of its neighbours; Euclid's algorithm leaves either the three-ray
fan or a four-ray Hirzebruch fan.  The displayed fans and self-intersections
for $0\leq n\leq6$ are checked by the same program; the determinant and
self-intersection formulas are valid for every $n$.

For products, a puncture loop of an affine hyperbolic curve has infinite order
in its free fundamental group.  At a crossing the two meridians belong to the
two different direct factors, so their $\Z^2$ injects.  Every stratum of the
product completion is a product of aspherical affine curves or points.
In the product of two hyperbolic curves, an original boundary component either
has positive genus or, when rational, meets at least three other components.
It can never be the first original component contracted in a common corner
refinement.  Hence every good completion is a corner blow-up of this product,
and the product is the unique minimal member.

For $\C^*\times B$, with $B$ the three-punctured line, the two sections have
valence three and cannot be contracted.  The projection to $\overline B$
extends through every corner modification.  Every modified missing fibre is
a chain between these sections; a reducible such fibre contains a bivalent
$(-1)$-component, since it is obtained by point blow-ups of a smooth fibre.
Thus in a minimal completion all three missing fibres are irreducible and the
surface is a $\PP^1$-bundle with two disjoint sections.  It is
$\PP(\mathcal O\oplus L)$; the fixed fibre coordinate on $U$ is exactly a
trivialization of $L|_B$.  A line bundle with this trivialization has a unique
meromorphic section equal to $1$ on $B$, hence is represented by an integral
divisor supported at the three missing points.  Conversely every such divisor
gives the required completion.  This explains why the marking remembers an
integer triple, even though the abstract ruled surface remembers only its
Hirzebruch degree up to sign.
\end{proof}

\begin{proposition}[\status{Refuted}: analytic and stacky overstatements]\label{P3a-scope}
Without an algebraic or bimeromorphic qualification, the assertion that any
two good marked completions are joined by stratal blow-ups and blow-downs is
false already for analytic surfaces.  With DM stacks admitted, ordinary
representable stratal blow-ups alone also fail to connect all good
compactifications, even for a torus.
\end{proposition}
\begin{proof}
Let $U=(\C^*)^2$ and use the toric completion $\PP^1\times\PP^1$ with two
markings related by the biholomorphism $(x,y)\mapsto(x,e^x y)$ of $U$.
Both pairs satisfy the criterion.  A sequence of analytic blow-ups and
blow-downs inducing the identity on the marked $U$ would give a meromorphic
map between the completions.  Restricting its second coordinate to the slice
$y=1$ would meromorphically extend $e^x$ across $x=\infty$, which is impossible
because that singularity is essential.  Thus the needed birational comparison
is absent, not merely difficult to factor.

For the stack statement take the same toric surface and form the order-two
root stack along one boundary component.  The open torus is unchanged; the
boundary strata are finite gerbes over tori and hence aspherical; local
meridian maps are the old injective maps restricted to finite-index
sublattices.  The rooted and unrooted pairs are both good.  Ordinary stratal
blow-ups do not introduce a root index along a generic boundary branch.
Locally their monomial matrices are unimodular and their stabilizer data are
inherited by pullback.  Their inverses have the same property.  They therefore
cannot change the unrooted integral normal structure into one having generic
$\mu_2$ inertia on that branch.  Root constructions must be added to the list
of operations if such stacky compactifications are to be compared.
\end{proof}

\subsection{Elementary fibrations and their profinite form}

\begin{theorem}[\status{Proved}: horizontal and vertical inertia]\label{P3b-discrete}
Let $f:U\to B$ be a complex algebraic elementary fibration in affine
hyperbolic curves, and suppose $B$ is an iterated elementary fibration of the
same kind over a point.  Fix a fibre $F$, base points and paths.  Then
\[
 1\longrightarrow\pi_1(F)\longrightarrow\pi_1(U)
 \longrightarrow\pi_1(B)\longrightarrow1                  \tag{EF}
\]
is exact and $U$ is aspherical.  Every horizontal puncture meridian injects
into $\pi_1(U)$.

Suppose also that the base has a good normal-crossings completion and the
family extends as a stable pointed log curve, with a smooth toroidal total
model.  At each vertical stratum its full boundary lattice injects.
The base hypothesis is essential: extension as a stable family by itself does
not imply vertical injectivity.  For an arbitrary such family for which (EF)
is exact, the exact criterion is that the kernel of the local lattice's map
to $\pi_1(B)$ inject into $\pi_1(F)$ under its lifted map in (EF).
\end{theorem}
\begin{proof}
The proper smooth compactification with finite \'etale marked boundary makes
$f$ a locally trivial differentiable bundle of punctured surfaces.  Each fibre
is aspherical, and induction up the tower makes the base aspherical.  Its
homotopy exact sequence proves (EF) and all higher homotopy vanishing.
A peripheral loop of a hyperbolic punctured surface is nontrivial and has
infinite order: its fundamental group is free, and the boundary word is
nontrivial (for one puncture and positive genus it is the nonempty product of
commutators).  Its inclusion in the fibre group and then in (EF) proves
horizontal injectivity.

At a smooth point of a vertical component the local lattice projects to the
base lattice.  At a horizontal intersection there is in addition the fibre
puncture meridian.  The base injectivity kills all coefficients projecting to
the base, leaving only a power of that puncture meridian.
At a nodal point of a semistable curve over a disk the equation is $xy=t$;
the map on lattices is $(a,b)\mapsto a+b$, with kernel generated by $(1,-1)$.
This kernel maps to the vanishing cycle in $\pi_1(F)$.  Stability makes that
cycle essential on the marked surface, so it has infinite order.
For several base parameters the nodal equation is toroidal,
$xy=\prod t_i^{a_i}$, after the prescribed base change and toroidal resolution.
On the torus lattice the kernel of the map to the base lattice has rank at
most one; it is the vanishing-cycle direction.  The corresponding statements
follow by the same lattice computation, and persist under regular
subdivision by Theorem~\ref{P2b}.  Stability excludes exactly the disk-bounding
vanishing cycles that would make this direction trivial.

For the final criterion, let $L$ be the local lattice, $q:L\to\pi_1(B)$ its
composite, and $\rho:L\to\pi_1(U)$ the local meridian map.  Every element of
$\ker\rho$ belongs to $\ker q$; exactness of (EF) identifies the restriction
$\rho|_{\ker q}$ with a map into $\pi_1(F)$.  This proves the equivalence.
To disprove sufficiency of stable extension without a good base, take
$U=F\times\C$ and the trivial family over $\PP^1$, marking all punctures of
$\overline F$ and deleting the fibre at infinity.  The family of compact
pointed fibres is stable and constant.  The vertical meridian at infinity is
zero in $\pi_1(U)=\pi_1(F)$, although every horizontal puncture meridian
injects.  This counterexample lies outside the hyperbolic tower hypothesis
on the base, exactly where that hypothesis was used.
\end{proof}

\begin{lemma}[\status{Proved}: the completion facts actually needed]\label{P3b-completion}
For the groups of an iterated elementary fibration in affine hyperbolic curves:
\begin{enumerate}[label=(\roman*)]
\item the groups are residually finite and good in Serre's sense;
\item profinite completion of (EF) is exact also on the left;
\item the procyclic completion of each puncture or essential simple-curve
subgroup of a fibre injects into its profinite completion.
\end{enumerate}
\end{lemma}
\begin{proof}
We give the extension argument to specify precisely what is being used.
If $1\to N\to G\to Q\to1$ has $N$ finitely generated and $Q$ good, the
profinite topology induced on $N$ is its full profinite topology.
Indeed, let $K$ be a characteristic finite-index subgroup of $N$; such
subgroups are cofinal because a finitely generated group has only finitely
many subgroups of any bounded index.  In $E=G/K$ the kernel $A=N/K$ is finite.
The centralizer $E_0=C_E(A)$ has finite index and fits into a central extension
\[
 1\to Z(A)\to E_0\to Q_0\to1
\]
with $Q_0$ of finite index in $Q$.  Goodness passes to finite-index subgroups
by the discrete and continuous Shapiro lemmas.  The extension class in
$H^2(Q_0,Z(A))$ therefore comes from continuous cohomology of $\widehat Q_0$.
A continuous cocycle with finite values factors through a finite quotient;
after normalization it restricts to zero on an open subgroup.  Consequently
the extension splits over some finite-index subgroup $Q_1\subset Q_0$.
The kernel of projection $Z(A)\times Q_1\to Z(A)$ is a finite-index subgroup
of $E$ meeting $A$ trivially.  Its normal core in $E$ has the same property.
The resulting finite quotient of $G$ induces an injection on $N/K$.
This proves fullness and hence injectivity of $\widehat N\to\widehat G$;
right exactness of completion gives the rest of the exact sequence.

A finitely generated free group is good: its discrete and profinite
cohomology with finite coefficients vanishes above degree one, and degree-one
cocycles correspond because their finite values and the finite module action
factor through finite quotients.  Free groups are residually finite, as is
seen by completing a finite labelled path in a rose to a finite graph cover.
Inductively apply the previous paragraph to the free fibre kernel.
The Lyndon--Hochschild--Serre spectral sequences for the discrete and completed
extensions, with finite coefficients, then agree: the fibre cohomology is
finite and concentrated in degrees zero and one, and the base is good.
Thus the total group is good.  Residual finiteness follows by detecting an
element either in the residually finite quotient or in the fully embedded
profinite completion of the residually finite kernel.  This is the free-fibre
version of the extension argument in
\citet[Lemma 10.9 and the proof of Lemma 10.11]{Bresciani2021}; that paper's
compact-orbicurve definition is not silently being substituted for the affine
fibres used here.

Finally, an essential simple curve in a finite-type orientable surface has a
lift in some finite cover whose homology class is nonzero and primitive.
For a nonseparating curve this is already true.  A puncture curve is nonzero
in homology if there is another puncture.  In the one-puncture positive-genus
case, a nontrivial double cover defined by a mod-two handle class has two
punctures, each mapping with degree one, which gives the assertion.  For a
separating nonperipheral essential curve with zero homology, take double
covers nontrivial on each of its two sides and glue along the two lifts of
the separating curve.  The resulting cover is connected and either lift is
nonseparating.  (If a side is planar, its at least two punctures provide the
required mod-two character.)  Projection of this primitive homology class
onto $\Z$ detects its reduction modulo every integer.  A finite-index
subgroup has its full profinite topology inside the whole group, so this
proves $\widehat\Z\hookrightarrow\widehat{\pi_1(F)}$ for the original
simple-curve subgroup.  If the chosen lift traverses the curve more than once,
its finite-index subgroup suffices, since multiplication by a positive integer
is injective in $\widehat\Z$.
\end{proof}

\begin{corollary}[\status{Proved}: the profinite form of $1^\circ$]\label{P3b-profinite}
Let an iterated elementary fibration as in Theorem~\ref{P3b-discrete} be defined
over a field $K$ finitely generated over $\Q$.  At horizontal strata, and at
vertical strata of the stable toroidal models with good base described there,
\[
 T^{I(x)}=\widehat\Z(1)^{I(x)}\longrightarrow\pi_1^{\mathrm{\acute et}}
 (U_{\overline K},\bar u)
\]
is injective.  This assertion uses discrete groups only through their verified
profinite embeddings.  No claim about wild inertia in residue characteristic
$2$ or $3$ is implied.
\end{corollary}
\begin{proof}
Choose an embedding $\overline K\hookrightarrow\C$, geometric base points,
and compatible roots of unity.  The comparison of finite \'etale covers with
finite topological covers identifies the geometric group with
$\widehat{\pi_1(U(\C))}$, and identifies $T$ with $\widehat\Z$ after that
choice.  Lemma~\ref{P3b-completion} embeds the completed fibre group and each
completed peripheral cyclic subgroup.  At a mixed or nodal stratum, project
a completed lattice relation to the completed base lattice.  Its remaining
kernel is the completion of the integral fibre direction computed in
Theorem~\ref{P3b-discrete}, and that procyclic direction injects by the lemma.
Integral unimodular changes and regular subdivisions are still invertible
over $\widehat\Z$.  This proves the assertion inductively up the tower.
Changing the paths conjugates the map; changing the compatible roots of unity
changes the chosen coordinate on the Tate module.  Neither change affects
injectivity.  Passing to a finite extension to label branches does not change
the geometric group.  The comparison used is the characteristic-zero
Riemann existence theorem for finite covers \citep[Expos\'e XII, Theorem 5.1 and Corollary 5.2]{SGA1}; it does not identify a discrete
group with its completion.
\end{proof}

\subsection{Higher dimension}

\begin{theorem}[\status{Known, with reduction}: the toric case]\label{P3c-toric}
For the torus $U=(\C^*)^n$, every smooth complete toric compactification
satisfies the criterion.  Any two such compactifications are joined by
blow-ups and blow-downs along smooth closures of torus strata, and every
intermediate toric pair satisfies the criterion.
\end{theorem}
\begin{proof}
For a smooth cone $\sigma$ with $r$ rays, those primitive rays form part of an
integral basis of $N=\pi_1(U)=\Z^n$.  The local meridian homomorphism is the
inclusion of their span.  Its orbit stratum is $(\C^*)^{n-r}$, hence
aspherical.  This proves the criterion on every cone of every smooth complete
fan.  W\l odarczyk's toric weak factorization theorem
\citep[Theorem 4.7.1]{Wlodarczyk2006} factors a toric birational map of smooth
complete toric varieties through smooth complete toric varieties by invariant
smooth-centre blow-ups and blow-downs.  In a smooth toric variety an irreducible
invariant centre occurring there is a smooth orbit closure, precisely an
allowed stratum closure.  The preceding cone calculation applies to each
intermediate fan.  Thus weak factorization is sufficient; no unproved strong
factorization assertion is required.
\end{proof}

\begin{conjecture}[\status{Open}: the remaining scheme comparison problem]\label{P3c-general}
Let $(X,D)$ and $(X',D')$ be good smooth proper \emph{algebraic variety}
compactifications of the same marked $U$ over $\C$.  Their identity birational
map on $U$ is a composite of blow-ups and blow-downs of smooth stratum closures.
After Lemma~\ref{P3a-core}, the first dimension not decided by the arguments
here is three.  For a fixed pair the necessary first obstruction is whether
the birational map is toroidal; a stratal factorization would force this.
\end{conjecture}
This formulation excludes the explicit analytic and stack counterexamples of
Proposition~\ref{P3a-scope}.  Ordinary weak factorization only provides centres
normal-crossing with the boundary, and Theorem~\ref{P2a} shows exactly why its
nonstratal centres cannot be retained as good intermediate models.  We also
tested a local Atiyah-flop construction: compactifying the exceptional
$\PP^1$ by two transverse boundary fibres leaves independent meridians on one
side, but the flopped side gives two dependent fibre meridians along the new
curve, so it is not a counterexample with two good ends.  Thus neither weak
factorization nor that local flop supplies a proof or refutation of the
three-dimensional statement.  The requested toric case is decided by
Theorem~\ref{P3c-toric}; the general statement is recorded as the exact
conjecture requested in Part 3(c), not as a proved theorem.
